\chapter{Una migración de principio a fin}
\label{ch:migracion}

Este capítulo recorre dos meses y medio de vida de Flexgrid. Se publican la 1.1.0 y la 2.0.0, se firma un
acuerdo, se retira una versión y se revoca. En cada paso se envían los mismos tres mensajes y se mira qué
cambia.

\noindent\begin{tabularx}{\linewidth}{@{}L{2.8cm}Y@{}}
\toprule
\textbf{Mensaje} & \textbf{Qué lleva y bajo qué versiones es válido} \\
\midrule
\cod{m1-forecast} & Previsión de demanda. Válido en las tres versiones. \\
\cod{m2-profiling} & Perfilado de consumidores. Válido en 1.0.0 y 1.1.0; inválido en 2.0.0. \\
\cod{m3-research} & Finalidad de investigación. Inválido en 1.0.0; válido desde 1.1.0. \\
\bottomrule
\end{tabularx}

\medskip
\cod{rows} envía los tres por \cod{dsp-ingress} e imprime, por cada uno, la línea de estado HTTP y la cabecera
\cod{VHub-Verdict}, seguidas del estado del cliente y sus versiones.

\section{El punto de partida}

3 de septiembre, con la 1.0.0 recién entrada en vigor:

\begin{lstlisting}[style=sh]
rows
\end{lstlisting}
\begin{lstlisting}[style=out]
m1-forecast    HTTP/1.1 200 OK | -
m2-profiling   HTTP/1.1 200 OK | -
m3-research    HTTP/1.1 400 Bad Request | VHub-Verdict: violation;gate=policy_intake;code="VERDICT_VIOLATION";rule="contract-request"
   Active [('1.0.0', 'head')]
\end{lstlisting}

\section{Una versión compatible}

El 20 de septiembre Lucía publica la 1.1.0, que añade la finalidad de investigación y marca obsoleto el
perfilado.

\begin{lstlisting}[style=sh]
syncat 2026-09-20T10:00:00Z
publish ../dsl/flexgrid-1.1.0.yaml
\end{lstlisting}
\begin{lstlisting}[style=out]
draft:     ffd85632051f7891
{"@context":"…","class":"compatible","findings":[],"valid":true}
candidate: 85089995ebc525b7
\end{lstlisting}

El hub la clasifica como \cod{compatible} por su número, y N4 ha comprobado contra la fuente de la 1.0.0 que lo
es de verdad. El \cod{diff} del candidato muestra el cambio en términos de ficheros: aparecen los seis de
\cod{profiles/odrl/metering/1.1.0/}, desaparecen los seis de \cod{1.0.0/} y cambian los cuatro del espacio y la
fuente.

Una versión compatible no admite ventana de gracia, porque no hay nada que hacer convivir:

\begin{lstlisting}[style=sh]
approve '{"grace_window":"P10D"}'
\end{lstlisting}
\begin{lstlisting}[style=out]
HTTP/2 422
{"…","title":"DATES_GRACE_ON_COMPATIBLE","status":422,
 "detail":"grace_window is only allowed on breaking versions","code":"DATES_GRACE_ON_COMPATIBLE","severity":"error"}
\end{lstlisting}

\begin{lstlisting}[style=sh]
approve '{}'
\end{lstlisting}
\begin{lstlisting}[style=out]
HTTP/2 201
{'version': '1.1.0', 'class': 'compatible', 'published_at': '2026-09-20T10:00:00Z',
 'new_content_from': '2026-09-22T10:00:00Z', 'enforced_from': '2026-09-22T10:00:00Z'}
\end{lstlisting}

Tres días después:

\begin{lstlisting}[style=sh]
syncat 2026-09-23T12:00:00Z
rows
\end{lstlisting}
\begin{lstlisting}[style=out]
m1-forecast    HTTP/1.1 200 OK | -
m2-profiling   HTTP/1.1 200 OK | -
m3-research    HTTP/1.1 200 OK | -
   Active [('1.1.0', 'head')]
\end{lstlisting}

La universidad ya puede negociar. El estado sigue siendo \cod{Active}, con una sola versión: la 1.1.0 ha
sustituido limpiamente a la 1.0.0, que pasa a \cod{Superseded}.

Y \cod{m2-profiling} pasa sin ninguna señal, aunque usa un término obsoleto. Es la trampa del
capítulo~\ref{ch:evolucion} vista en producción: el aviso existe en \cod{shapes.ttl}, y nadie lo ejecuta.

\section{Una versión de ruptura}

El 1 de octubre se publica la 2.0.0, que retira el perfilado.

\begin{lstlisting}[style=sh]
syncat 2026-10-01T10:00:00Z
publish ../dsl/flexgrid-2.0.0.yaml
\end{lstlisting}
\begin{lstlisting}[style=out]
draft:     a5fad501c89c41e8
{"@context":"…","class":"breaking","findings":[],"valid":true}
candidate: 5552a13ebfe21682
\end{lstlisting}

Lucía intenta acortar la convivencia a tres días:

\begin{lstlisting}[style=sh]
approve '{"grace_window":"P3D"}'
\end{lstlisting}
\begin{lstlisting}[style=out]
HTTP/2 422
{"…","title":"DATES_GRACE_TOO_SHORT","status":422,
 "detail":"grace_window P3D is below min_grace_window P7D","code":"DATES_GRACE_TOO_SHORT","severity":"error"}
\end{lstlisting}

El mínimo son siete días y es del hub: ninguna \emph{steward} puede dejar a los participantes sin tiempo para
migrar. Con la gracia por defecto, de treinta:

\begin{lstlisting}[style=sh]
approve '{}'
\end{lstlisting}
\begin{lstlisting}[style=out]
HTTP/2 201
{'version': '2.0.0', 'class': 'breaking', 'published_at': '2026-10-01T10:00:00Z',
 'new_content_from': '2026-10-03T10:00:00Z', 'enforced_from': '2026-11-02T10:00:00Z'}
\end{lstlisting}

Ahora las dos fechas son distintas, y entre ellas está la ventana de gracia.

\begin{center}
\begin{tikzpicture}[font=\scriptsize, x=0.156cm, y=1cm]
  \draw[alxBorder] (0,-0.5) -- (80,-0.5);
  \foreach \d/\l in {2/3 sep, 21/22 sep, 32/3 oct, 62/2 nov, 72/12 nov}
    { \draw[alxBorder] (\d,-0.5) -- (\d,-0.65); \node[below] at (\d,-0.68) {\l}; }
  \fill[alxAmber, opacity=0.13] (32,0.25) rectangle (62,1.75);
  \node[alxAmber!80!black] at (47,1.0) {ventana de gracia: conviven};
  \draw[alxBlue, dashed, line width=1.2pt] (0,2.2) -- (2,2.2);
  \draw[alxBlue, line width=3.2pt] (2,2.2) -- (21,2.2);
  \node[left] at (0,2.2) {1.0.0};
  \draw[alxTeal, dashed, line width=1.2pt] (19,1.4) -- (21,1.4);
  \draw[alxTeal, line width=3.2pt] (21,1.4) -- (62,1.4);
  \draw[alxTeal!45, line width=3.2pt] (62,1.4) -- (72,1.4);
  \node[left] at (19,1.4) {1.1.0};
  \node[above, alxTeal!70!black] at (67,1.45) {retenida};
  \node[alxRed, right] at (72,1.4) {revocada};
  \draw[alxPurple, dashed, line width=1.2pt] (30,0.6) -- (32,0.6);
  \draw[alxPurple, line width=3.2pt] (32,0.6) -- (80,0.6);
  \node[left] at (30,0.6) {2.0.0};
\end{tikzpicture}

{\scriptsize El tramo discontinuo es la ventana de propagación: publicada, pero aún no vigente.}
\end{center}

\section{Lo que ve el borde, instante a instante}

\begin{lstlisting}[style=sh]
syncat 2026-10-02T12:00:00Z; rows
syncat 2026-10-03T06:00:00Z; clock 2026-10-03T09:59:55Z; rows
syncat 2026-10-10T12:00:00Z; rows
syncat 2026-11-02T06:00:00Z; clock 2026-11-02T10:00:05Z; rows
syncat 2026-11-05T12:00:00Z; rows
\end{lstlisting}

\begin{longtable}{@{}L{2.7cm}C{1.7cm}C{1.7cm}C{1.7cm}L{3.3cm}L{3.2cm}@{}}
\toprule
\textbf{Instante} & \textbf{m1} & \textbf{m2} & \textbf{m3} & \textbf{Estado y versiones} & \textbf{Por qué} \\
\midrule
\endhead
3 sep, 12:00 & \ok & \ok & \viol & \cod{Active}: 1.0.0 & Solo la 1.0.0. \\
23 sep, 12:00 & \ok & \ok & \ok & \cod{Active}: 1.1.0 & La 1.1.0 la ha sustituido. \\
2 oct, 12:00 & \ok & \ok & \ok & \cod{Active}: 1.1.0; 2.0.0 programada & La 2.0.0 se conoce, pero no vale aún. \\
3 oct, 09:59:55 & \ok & \warn & \ok & \cod{Transition}: 2.0.0 cabeza, 1.1.0 vigente & La 2.0.0 entra un margen de arrendamiento antes. \\
10 oct, 12:00 & \ok & \warn & \ok & \cod{Transition} & Ventana de gracia. \\
2 nov, 10:00:05 & \ok & \warn & \ok & \cod{Transition} & Aún dentro del margen tras \cod{enforced_from}. \\
5 nov, 12:00 & \ok & \viol & \ok & \cod{Active}: 2.0.0; 1.1.0 retenida & Solo la 2.0.0. \\
\bottomrule
\end{longtable}

Es el patrón \emph{expandir, migrar, contraer}. Primero se acepta lo nuevo. Lo viejo sigue funcionando con un
aviso. El rechazo empieza solo después de \cod{enforced_from}.

\subsection{El aviso}

Durante la ventana, \cod{m2-profiling} se reenvía al conector y la respuesta vuelve con una cabecera:

\begin{lstlisting}[style=out]
m2-profiling   HTTP/1.1 200 OK | VHub-Verdict: warning;gate=policy_intake;code="VERDICT_NOT_IN_HEAD";rule="contract-request";version="1.1.0";sunset=@1793613600
\end{lstlisting}

\noindent\begin{tabularx}{\linewidth}{@{}L{5.7cm}Y@{}}
\toprule
\textbf{Parámetro} & \textbf{Significado} \\
\midrule
\cod{warning} & La clasificación. \\
\cod{gate=policy_intake} & El gate que la produjo. \\
\texttt{code="VERDICT\_NOT\_IN\_HEAD"} & El mensaje no conforma con la versión cabeza, pero sí con una anterior aún vigente. \\
\texttt{version="1.1.0"} & La versión que lo aceptó. \\
\cod{sunset=@1793613600} & Instante en que dejará de aceptarse: 2 de noviembre de 2026 a las 10:00 UTC. \\
\bottomrule
\end{tabularx}

\medskip
La comercializadora no tiene que leer ningún boletín. Su propio tráfico le dice que lo que envía tiene fecha de
caducidad, y cuál.

\cod{:explain} muestra el razonamiento completo:

\begin{lstlisting}[style=sh]
ex m2-profiling | short
\end{lstlisting}
\begin{lstlisting}[style=out]
rule=contract-request state=Transition action=forward-notice code=None
  gate=policy_intake  class=warning   codes=['VERDICT_NOT_IN_HEAD'] matched=1.1.0
     - Violation json-schema /permission/0/action validation failed
     - Info loader  POC_EVAL_UNSUPPORTED_ROLE: shacl artefact profiles/odrl/metering/2.0.0/shapes.ttl is not evaluated
  verdict: warning;gate=policy_intake;code="VERDICT_NOT_IN_HEAD";rule="contract-request";version="1.1.0";sunset=@1793613600
\end{lstlisting}

Los hallazgos son los de la cabeza, la 2.0.0, que rechaza la acción. \cod{matched=1.1.0} dice quién lo aceptó.

\begin{bajoElCapo}[El algoritmo de clasificación]
Para un gate, después de las filas de estado:
\begin{enumerate}[leftmargin=*, itemsep=1pt]
  \item Cuerpo ilegible, o falta el nodo de \cod{select}: violación, \cod{MESSAGE_UNPARSEABLE}.
  \item Ninguna versión vigente: violación, \cod{NO_VERSION_IN_FORCE}.
  \item Se evalúa contra la \textbf{cabeza}. Si conforma: ok.
  \item Si no, contra las anteriores aún vigentes, de la más nueva a la más vieja. Si una lo acepta: aviso,
        \cod{VERDICT_NOT_IN_HEAD}.
  \item Si ninguna: violación. \cod{CONTEXT_NOT_IN_PACKAGE} si el motivo es un contexto ajeno;
        \cod{VERDICT_VIOLATION} en otro caso.
\end{enumerate}
\end{bajoElCapo}

\subsection{El margen de arrendamiento}

Dos filas de la tabla tienen horas extrañas. El 3 de octubre a las 09:59:55 la 2.0.0 ya acepta mensajes, cinco
segundos antes de su \cod{new_content_from}. El 2 de noviembre a las 10:00:05, cinco segundos después de
\cod{enforced_from}, la 1.1.0 todavía se acepta.

Es el \cod{lease_margin}, de diez segundos: el cliente ensancha cada frontera por los dos lados. Los relojes de
dos participantes nunca coinciden del todo, y sin margen un mensaje enviado en el segundo exacto del cambio
podría ser válido para quien lo envía e inválido para quien lo recibe.

\section{Los acuerdos sobreviven a su versión}

El 10 de octubre, en plena ventana, la comercializadora y la distribuidora cierran un acuerdo que incluye el
perfilado. Es \cod{examples/msgs/a1-agreement.json}:

\begin{lstlisting}[style=json]
{
  "@context": [ … ],
  "@type": "ContractAgreementMessage",
  "providerPid": "urn:uuid:9c1f7a52-dso-neg-0042",
  "consumerPid": "urn:uuid:7d2a9f10-ret-neg-0007",
  "agreement": {
    "@id": "urn:uuid:agreement-retail-0007", "@type": "Agreement",
    "profile": "https://flexgrid.example/odrl/profile/metering",
    "target": "urn:uuid:dataset-qh-load-madrid-2026",
    "assigner": "did:web:dso.flexgrid.example",
    "assignee": "did:web:retailer.flexgrid.example",
    "permission": [ { "action": "profileConsumer", "constraint": [
        {"leftOperand": "purpose", "operator": "eq", "rightOperand": {"@id": "enr:billing"}},
        {"leftOperand": "retentionDays", "operator": "lteq", "rightOperand": 30} ] } ]
  }
}
\end{lstlisting}

\begin{lstlisting}[style=sh]
post /negotiations/urn:uuid:7d2a9f10-ret-neg-0007/agreement a1-agreement
curl -s "$ADMINAPI/agreements/urn:uuid:agreement-retail-0007"
\end{lstlisting}
\begin{lstlisting}[style=out]
a1-agreement   HTTP/1.1 200 OK | VHub-Verdict: warning;gate=policy_intake;code="VERDICT_NOT_IN_HEAD";rule="contract-agreement-in";version="1.1.0";sunset=@1793613600

{ "agreement_id": "urn:uuid:agreement-retail-0007",
  "provider_pid": "urn:uuid:9c1f7a52-dso-neg-0042",
  "consumer_pid": "urn:uuid:7d2a9f10-ret-neg-0007",
  "role": "consumer",
  "version": "1.1.0",
  "manifest": "ni:///sha-256;8kvT3m-30MLD6JA7VFtn-iNDqtFe2fYvyx_uG5ImLgI",
  "declared_profile": "https://flexgrid.example/odrl/profile/metering",
  "recorded_at": "2026-10-10T12:00:00Z",
  "source": "dsp-ingress",
  "usable": true }
\end{lstlisting}

El acuerdo pasa con aviso, y el cliente lo \textbf{captura}: guarda su identificador, los de las dos partes, su
papel y, sobre todo, la versión que lo validó, la 1.1.0. La captura ocurre cuando el conector responde \cod{2xx}
y antes de devolver la respuesta.

\begin{ideaClave}
Un acuerdo queda ligado a la versión del vocabulario bajo la que se firmó, no a la vigente. Un contrato no se
reinterpreta cada vez que cambian las reglas.
\end{ideaClave}

Con el acuerdo capturado, una transferencia bajo él se admite. \cod{examples/msgs/t1-transfer.json} es un
\cod{TransferRequestMessage} con \cod{agreementId}:

\begin{lstlisting}[style=sh]
post /transfers/request t1-transfer
post /transfers/request t2-transfer-no-format -b
\end{lstlisting}
\begin{lstlisting}[style=out]
t1-transfer             HTTP/1.1 200 OK | -
t2-transfer-no-format   HTTP/1.1 400 Bad Request | VHub-Verdict: violation;gate=dsp_message;code="VERDICT_VIOLATION";rule="transfer-request"
    {"@context":["https://w3id.org/dspace/2025/1/context.jsonld"],"@type":"TransferError",
     "code":"VERDICT_VIOLATION","consumerPid":"urn:uuid:7d2a9f10-ret-trf-0001","providerPid":null,
     "reason":["'allOf' failed","missing property 'format'"]}
\end{lstlisting}

La segunda transferencia es la primera vez que el slot \cod{dsp} vendorizado rechaza algo: al mensaje le falta
\cod{format}, que el esquema de \cod{vendor/dsp/2025-1/} exige. El error es ahora un \cod{TransferError}, el que
DSP define para ese intercambio.

\subsection{Después de la retirada}

El 5 de noviembre la 1.1.0 ya no está vigente. Una petición nueva con perfilado se rechaza. Pero:

\begin{lstlisting}[style=sh]
curl -s $ADMINAPI/versions | python3 -c "import json,sys;[print(v['version'], v['class'], v['state'], 'retained' if v['retained'] else '') for v in json.load(sys.stdin)]"
post /transfers/request t1-transfer
\end{lstlisting}
\begin{lstlisting}[style=out]
1.0.0 initial Superseded
1.1.0 compatible Retired retained
2.0.0 breaking Published retained

t1-transfer             HTTP/1.1 200 OK | -
\end{lstlisting}

La 1.1.0 está \cod{Retired}, pero su paquete se \textbf{retiene} porque el acuerdo la referencia, y la
transferencia se sigue admitiendo. La 1.0.0, a la que no apunta ningún acuerdo, ya no se conserva.

\noindent\begin{tabularx}{\linewidth}{@{}L{2.6cm}Y@{}}
\toprule
\textbf{Estado} & \textbf{Significado, en un instante de referencia} \\
\midrule
\cod{Scheduled} & Antes de \cod{new_content_from}. \\
\cod{Published} & Vigente y cabeza. \\
\cod{Deprecated} & Vigente, pero no cabeza. \\
\cod{Superseded}, \cod{Retired} & Ya no vigente. \\
\cod{Revoked} & Revocada. \\
\bottomrule
\end{tabularx}

\section{Revocar}

El 10 de noviembre el regulador declara ilícito el perfilado individual. Retirar ya no basta: los acuerdos
firmados bajo la 1.1.0 no deben amparar más transferencias.

\begin{lstlisting}[style=sh]
syncat 2026-11-10T00:00:00Z
R=$HUB/v1/namespaces/flexgrid/versions
curl -sk -X POST "$R/1.1.0:revoke" -H "Authorization: Bearer $(cat steward.token)" -d '{}'
curl -sk -X POST "$R/2.0.0:revoke" -H "Authorization: Bearer $(cat steward.token)" -d '{"justification":"test"}'
\end{lstlisting}
\begin{lstlisting}[style=out]
{"…","title":"DATES_JUSTIFICATION_REQUIRED","status":422,"detail":"a revocation requires a justification", … }
{"…","title":"DATES_REVOKE_LEAVES_NO_HEAD","status":422,"detail":"no version would be in force at 2026-11-12T00:00:00Z", … }
\end{lstlisting}

Los dos intentos fallan, y con razón: una revocación exige justificación, y no se puede revocar la única
versión vigente, porque el espacio se quedaría sin reglas.

\noindent\begin{tabularx}{\linewidth}{@{}C{1.3cm}Y L{6.1cm}@{}}
\toprule
\textbf{Regla} & \textbf{Comprobación} & \textbf{Código} \\
\midrule
REV-1 & La versión no está ya revocada. & \cod{DATES_ALREADY_REVOKED} \\
REV-2 & Hay justificación. & \cod{DATES_JUSTIFICATION_REQUIRED} \\
REV-3 & \cod{revoked_from} es al menos el instante de la petición más la ventana de propagación. & \cod{DATES_REVOCATION_TOO_EARLY} \\
REV-4 & Queda alguna versión vigente. & \cod{DATES_REVOKE_LEAVES_NO_HEAD} \\
\bottomrule
\end{tabularx}

\begin{lstlisting}[style=sh]
curl -sk -X POST "$R/1.1.0:revoke" -H "Authorization: Bearer $(cat steward.token)" \
  -d '{"justification":"Individual consumer profiling was ruled unlawful by the regulator"}'
\end{lstlisting}
\begin{lstlisting}[style=out]
{ … "credentialSubject": { "version": "1.1.0", "class": "compatible",
    "enforced_from": "2026-09-22T10:00:00Z", "revoked_from": "2026-11-12T00:00:00Z", … } }
\end{lstlisting}

La revocación tampoco es instantánea: \cod{revoked_from} es el 12 de noviembre, dos días después. Las fechas ya
emitidas no cambian; \cod{revoked_from} se añade una sola vez.

El listado siguiente lleva la revocación, y el cliente, al verla, avisa de cada acuerdo afectado:

\begin{lstlisting}[style=sh]
grep AGREEMENT var/client/events.jsonl
\end{lstlisting}
\begin{lstlisting}[style=json]
{"type":"vhub.agreement","ts":"…","code":"AGREEMENT_REVOCATION_SCHEDULED","severity":"info",
 "detail":"revoked_from 2026-11-12T00:00:00Z","version":"1.1.0",
 "agreement":"urn:uuid:agreement-retail-0007"}
\end{lstlisting}

Ese evento es lo que el operador de la comercializadora debe conectar a su sistema de alertas: tiene dos días
para renegociar.

\begin{lstlisting}[style=sh]
clock 2026-11-12T00:00:05Z; sleep 2; post /transfers/request t1-transfer
clock 2026-11-12T00:00:10Z; sleep 2; post /transfers/request t1-transfer -b
\end{lstlisting}
\begin{lstlisting}[style=out]
t1-transfer   HTTP/1.1 200 OK | -

t1-transfer   HTTP/1.1 400 Bad Request | VHub-Verdict: violation;gate=agreement_admission;code="AGREEMENT_VERSION_REVOKED";rule="transfer-request"
    {"@context":["https://w3id.org/dspace/2025/1/context.jsonld"],"@type":"TransferError",
     "code":"AGREEMENT_VERSION_REVOKED","consumerPid":"urn:uuid:7d2a9f10-ret-trf-0001","providerPid":null,
     "reason":["AGREEMENT_VERSION_REVOKED"]}
\end{lstlisting}

A las 00:00:05 la transferencia aún pasa, dentro del margen de arrendamiento. A las 00:00:10 se rechaza, y el
acuerdo figura como \texttt{"usable": false}.

\begin{ideaClave}
\textbf{Retirar} no rompe los acuerdos existentes. \textbf{Revocar} sí, y se anuncia con antelación. Son dos
herramientas distintas: la primera es para evolucionar el vocabulario, la segunda para cuando algo que se
permitió no debió permitirse.
\end{ideaClave}

\section{El rastro en el hub}

Los eventos del \emph{namespace} cuentan la misma historia desde el otro lado. Resumidos:

\begin{lstlisting}[style=out]
10 2026-09-20T10:00:00Z draft.created        lucia
11 2026-09-20T10:00:00Z candidate.proposed   lucia
12 2026-09-20T10:00:00Z release.aborted      lucia     <- el intento con grace_window
13 2026-09-20T10:00:00Z candidate.approved   lucia
14 2026-09-20T10:00:00Z release.published    lucia     versions/1.1.0
15 2026-09-20T10:00:00Z listing.issued       vhub
…
20 2026-10-01T10:00:00Z release.aborted      lucia     <- el intento con P3D
21 2026-10-01T10:00:00Z candidate.approved   lucia
22 2026-10-01T10:00:00Z release.published    lucia     versions/2.0.0
…
30 2026-11-10T00:00:00Z version.revoked      lucia     versions/1.1.0
31 2026-11-10T00:00:00Z listing.issued       vhub
\end{lstlisting}

Los dos intentos de aprobación fallidos también quedaron registrados, como \cod{release.aborted}. Una aprobación
que falla no deja nada visible para los clientes, pero sí deja constancia de que se intentó.

\begin{retos}
\begin{enumerate}[leftmargin=*, itemsep=2pt]
  \item Repite la migración aprobando la 2.0.0 con \texttt{\{"grace\_window":"P10D"\}} y calcula antes a mano el
        nuevo \cod{enforced_from} y el \cod{sunset} de la cabecera.
  \item Pon \texttt{unknown\_agreement: reject} en el cliente y envía una transferencia con un \cod{agreementId} que
        nunca se capturó.
  \item ¿Qué pasaría con el acuerdo si en lugar de revocar la 1.1.0 se publicara una 3.0.0? Compruébalo.
\end{enumerate}
\end{retos}
